\documentclass[10pt,a4paper]{amsart}
\usepackage[margin=19mm]{geometry}
\usepackage{amsmath,amssymb,mathtools}
\usepackage[scr=rsfs,cal=euler]{mathalfa}
\usepackage{xfrac}
\usepackage[unicode,hidelinks,bookmarks=false]{hyperref}
\newcommand{\ie}{i.e.\ }
\newcommand{\cf}{cf.\ }
\newcommand{\resp}{respectively\ }
\newcommand{\Subsection}{\subsection}
\providecommand{\nobreakdash}{\nobreak\hbox{-}}
\setlength{\emergencystretch}{2em}
\multlinegap0pt
\usepackage{amssymb}
\newtheorem{lemma}{Lemma}
\newtheorem{theorem}{Theorem}
\title{Combinatorial characterzations of $T$-designs in the nonbinary Johnson scheme}
\author{Hiroshi Nozaki, Yuta Watanabe}
\date{Source: arXiv:2512.22034v1 (CC BY 4.0)}
\begin{document}
\maketitle

\noindent\emph{Source and licence.} Combinatorial characterzations of $T$-designs in the nonbinary Johnson scheme. Version arXiv:2512.22034v1 (CC BY 4.0), distributed by arXiv under the Creative Commons Attribution 4.0 International licence, http://creativecommons.org/licenses/by/4.0/, as recorded in the arXiv licence metadata for this exact version and shown on the abstract page https://arxiv.org/abs/2512.22034v1. Full licence notice: \texttt{licenses/arxiv-2512.22034-CC-BY-4.0.txt}.\par\smallskip

\noindent\textit{Editorial scope.} Counted unit: the article's theorem thm:T-design (source0.tex lines 259--268). The article's complete original proof of it is delivered with it (source0.tex lines 269--394, one \texttt{proof} environment of 126 lines). No mathematics is written, repaired or re-derived: the statement and the proof are the article's own lines, in the article's own notation, and no retained line is substituted or re-flowed.  Retained uncounted as the source-local context the proof needs: lines 220--226; lines 241--243.  Nothing was omitted from any retained span, so the package carries no omission record.  No retained line contains a citation, so the wrapper carries no bibliography and no bibliography entry.  No retained line contains a figure environment, an image-inclusion command or drawing code, so no image is delivered and the package declares no asset dependency.  Editorial formatting only, in addition: an AMS \texttt{amsart} preamble that restates the article's own declarations and macros used by the retained lines, namely the environments \texttt{lemma}, \texttt{theorem} on separate counters, together with the standard packages those lines need.  The retained environments print in this document as Lemma 1, Theorem 1 in order of appearance, whereas the article prints the same items with the numbering of its own full document.  The complete native source of the article as distributed in the arXiv e-print for this exact version is bundled unchanged as \texttt{sources/191705/source0.tex}.
\begin{lemma} \label{lem:component}
Let $(r,s) \in L$. 
The set of components of $C_{rs} \in \mathfrak{B}$ is 
\[
\mathfrak{C}(C_{rs})=\{E_{is}  \mid (i,s) \in L, i\leq r  \}.
\]
\end{lemma}

\begin{equation} \label{eq:E_ij}
E_{ij}\phi_Y=\frac{|Y|}{|X|}E_{ij}\phi_X    
\end{equation}

\begin{theorem}\label{thm:T-design}
Suppose $(0,0)\ne (r,s) \in L$ with $r \leq m=\min\{w,n-w\}$. 
Let 
\begin{align*}
T&=\bigsqcup_{s'\in \{0,1,\ldots,s\}}\{(i,s') \in L \mid  i \leq r\}\setminus \{(0,0)\}\\    
&= \{(i,j) \in L \mid i \leq r, j\leq s\} \setminus \{(0,0)\}. 
\end{align*}
Then, $Y$ is a $T$-design in $J_q(w,n)$ if and only if 
$m_{\mathcal{R},\mathcal{S}}(Y,\omega)$ is constant for all $\omega \in \overline{F}^s$, $\mathcal{R}$, and $\mathcal{S}$ with $\mathcal{S}\subset \mathcal{R}$, $|\mathcal{R}|=r$, and $|\mathcal{S}|=s$.  
\end{theorem} 

\begin{proof}
From Lemma~\ref{lem:component} and \eqref{eq:E_ij}, 
$Y$ is a $T$-design if and only if for each $s' \in \{0,1,\ldots,s\}$, 
\begin{equation} \label{eq:crs}
C_{rs'}\phi_Y=\frac{|Y|}{|X|} C_{rs'} \phi_X.     
\end{equation} 
Note that $(r,s') \in L$ for each $s' \in \{0,1,\ldots, s\}$ with our assumption $r\leq m$. 
    The following are equivalent conditions of \eqref{eq:crs}: 
    \begin{align*}
C_{rs'}(\phi_Y-\frac{|Y|}{|X|} \phi_X)=0 
&\Leftrightarrow A_{rs'}\overline{A_{rs'}^\top}(\phi_Y-\frac{|Y|}{|X|} \phi_X) =0\\
& \Leftrightarrow \overline{A_{rs'}^\top}(\phi_Y-\frac{|Y|}{|X|} \phi_X)=0\\
& \Leftrightarrow A_{rs'}^\top(\phi_Y-\frac{|Y|}{|X|} \phi_X)=0
\\
& \Leftrightarrow A_{rs'}^\top\phi_Y=\frac{|Y|}{|X|}A_{rs'}^\top \phi_X. 
    \end{align*}
 
    For $a \in W_{rs}$, we have $w_H(a)=|\overline{a}|=r$ and $w_M(a)=|\tilde{a}|=s$, where $\tilde{a}=\{i \mid a_i\ne 0,1\}$. 
    Let $\tilde{a}=\{i_1,\ldots, i_s\}$. 
    The $a$-th entry of $A_{rs}^\top \phi_X$ is calculated as follows: 
    \begin{align}
    (A_{rs}^\top \phi_X)(a)&=\sum_{x \in X} (a,x) \nonumber \\
    &=\sum_{x \in X} \prod_{i \in \overline{a}} 
    \psi_{a_i}(x_i)  \nonumber \\
    &= \sum_{x \in X: \overline{a} \subset \overline{x} } \prod_{i \in \overline{a}} \psi_{a_i}(x_i) \label{eq:5} \\
    &= \sum_{x \in X: \overline{a} \subset \overline{x} } \prod_{j =1}^s \psi_{a_{i_j}}(x_{i_j})  \nonumber \\
    &=(q-1)^{w-s} \binom{n-r}{w-r}\sum_{\omega_1,\ldots,\omega_s \in \overline{F}} \prod_{j =1}^s \psi_{a_{i_j}}(\omega_j) \label{eq:4} \\
    &= (q-1)^{w-s} \binom{n-r}{w-r} \prod_{j=1}^s \sum_{\omega \in \overline{F}}\psi_{a_{i_j}}(\omega)  \nonumber \\
    &=\begin{cases}
        (q-1)^{w} \binom{n-r}{w-r} &\text{if $s=0$},  \nonumber\\
        0 &\text{if $s\geq 1$ }. 
    \end{cases}
    \end{align} 
In the equality \eqref{eq:4}, we count the number of terms that take the same value, namely those in  
\[
\{ x \in X \mid \overline{a} \subset \overline{x}, x_{i_j} = \omega_j( j =1,\ldots,s) \}
\]
for given $\omega_i \in \overline{F}$. 
This is equal to 
\[
m_{\mathcal{R},\mathcal{S}}(X,\omega)=(q-1)^{w-s} \binom{n-r}{w-r}, 
\]
where $\mathcal{S} \subset \mathcal{R}$ with $|\mathcal{R}|=r$ and $|\mathcal{S}|=s$. 

(I) We first prove sufficiency.  
Fix $\mathcal{R}, \mathcal{S} \subset \{1,2,\ldots,n\}$ with $\mathcal{S} \subset \mathcal{R}$, $|\mathcal{R}| = r$, and $|\mathcal{S}| = s$.  
From the equality \eqref{eq:5}, for $a \in W_{rs'}$ ($0 \le s' \le s$) with $\overline{a} = \mathcal{R}$ and $\tilde{a} \subset \mathcal{S}=\{i_1,\ldots, i_s\}$, we have
\begin{align*}
    (A_{rs'}^\top \phi_Y)(a)
    &= \sum_{y \in Y:\, \overline{a} \subset \overline{y}} \prod_{i \in \overline{a}} \psi_{a_i}(y_i)\\
    &= \sum_{y \in Y:\, \mathcal{R} \subset \overline{y}} \prod_{j =1}^s \psi_{a_{i_j}}(y_{i_j})\\
    &= \sum_{\omega \in \overline{F}^{s}} m_{\mathcal{R},\mathcal{S}}(Y,\omega) 
       \prod_{j=1}^s \psi_{a_{i_j}}(\omega_j).
\end{align*}
Since $Y$ is a $T$-design, we have
\[
\sum_{\omega \in \overline{F}^{s}} 
   m_{\mathcal{R},\mathcal{S}}(Y,\omega) 
   \prod_{j=1}^s \psi_{a_{i_j}}(\omega_j)
   = (A_{rs'}^\top \phi_Y)(a)
   = \frac{|Y|}{|X|}(A_{rs'}^\top \phi_X)(a)
   = 
   \begin{cases}
      \dfrac{|Y|}{|X|}(q-1)^{w} \binom{n-r}{w-r}, & \text{if } s' = 0,\\[6pt]
      0, & \text{if } s' \ge 1.
   \end{cases}
\]
The number of these equalities is $|\overline{F}^{s}|$, corresponding to the choices of $a_{i_j} \in \overline{F}$ ($j=1,\ldots,s$).  
For the linear equations in the indeterminates $m_{\mathcal{R},\mathcal{S}}(Y,\omega)$ ($\omega \in \overline{F}^s$),  
the coefficient matrix 
\[
\left( \prod_{j=1}^s \psi_{a_{i_j}}(\omega_j) \right)_{(a_{i_1},\ldots,a_{i_s}) \in \overline{F}^s,\, \omega \in \overline{F}^s}
\]
is nonsingular.  
Therefore, the solution is unique, and it is indeed determined by
\[
m_{\mathcal{R},\mathcal{S}}(Y,\omega)
   = \frac{|Y|}{|X|}(q-1)^{w-s} \binom{n-r}{w-r},
\]
as can be seen by comparing with $m_{\mathcal{R},\mathcal{S}}(X,\omega)$.  
This completes the proof of sufficiency.


(II) We prove necessity. 
Assume that $m_{\mathcal{R},\mathcal{S}}(Y,\omega)$ takes a constant value $\lambda_{rs}$. 
Let $\mathcal{S}=\{i_1,\ldots, i_s\}\subset \mathcal{R}$. We have 
\[
 |\{y\in Y \mid  \mathcal{R}\subset \overline{y} \}|
 = \sum_{\omega \in \overline{F}^s}|\{y \in Y \mid  \mathcal{R} \subset \overline{y}, y_{i_j}=\omega_j\ (j=1,\ldots, s)\}|=(q-1)^s \lambda_{rs}. 
\]
By double counting the set $\{(y,\mathcal{R}) \mid  y \in Y, |\mathcal{R}|=r, \mathcal{R} \subset \overline{y}\}$, we obtain 
\[
|Y|\binom{w}{r}
=\sum_{y\in Y} |\{\mathcal{R} \mid \mathcal{R}\subset \overline{y}, |\mathcal{R}|=r  \}|
=\sum_{|\mathcal{R}|=r} |\{y\in Y \mid  \mathcal{R}\subset \overline{y} \}|
=\binom{n}{r}(q-1)^s\lambda_{rs}. 
\]
Together with $|X|=\binom{n}{w}(q-1)^w$, it follows that
\[
\lambda_{rs}=|Y|\frac{\binom{w}{r}}{\binom{n}{r}}(q-1)^{-s}
=\frac{|Y|}{|X|}\frac{\binom{n}{w}\binom{w}{r}}{\binom{n}{r}}(q-1)^{w-s}
=\frac{|Y|}{|X|}\binom{n-r}{w-r}(q-1)^{w-s}. 
\]

Moreover, $m_{\mathcal{R},\mathcal{S}'}(Y,\omega')$ is also constant for $\mathcal{S}'\subset \mathcal{R}$ with $|\mathcal{S}'|=s'\leq s$. 
Indeed, for $\mathcal{S}'=\{i_1,\ldots,i_{s'}\} \subset \mathcal{S}=\{i_1,\ldots, i_s\} \subset \mathcal{R}$, 
\begin{align}
m_{\mathcal{R},\mathcal{S'}}(Y,\omega')&=|\{y \in Y \mid \mathcal{R} \subset \overline{y},y_{i_j}=\omega_{j}'\, (j=1,\ldots,s')\}|\nonumber\\
&=\sum_{\omega \in \overline{F}^s\colon\,  \omega_{j}=\omega_{j}'\, (j=1,\ldots,s')} 
|\{y \in Y \mid \mathcal{R} \subset \overline{y},y_{i_j}=\omega_{j}\, (j=1,\ldots,s)\}|\nonumber\\
&=(q-1)^{s-s'} \lambda_{r,s} \label{eq:(q-1)ss'lambda}\\
&=\frac{|Y|}{|X|} (q-1)^{w-s'}\binom{n-r}{w-r}\nonumber\\
&=\frac{|Y|}{|X|} m_{\mathcal{R},\mathcal{S'}}(X,\omega').\nonumber 
\end{align}

It therefore follows that 
\begin{multline*}
(A_{rs'}^\top \phi_Y)(a)=\sum_{\omega' \in \overline{F}^{s'}} 
   m_{\mathcal{R},\mathcal{S'}}(Y,\omega') 
   \prod_{j=1}^{s'} \psi_{a_{i_j}}(\omega_j')\\
   =\sum_{\omega' \in \overline{F}^{s'}} 
   \frac{|Y|}{|X|}m_{\mathcal{R},\mathcal{S}'}(X,\omega') 
   \prod_{j=1}^{s'} \psi_{a_{i_j}}(\omega_j')=\frac{|Y|}{|X|}(A_{rs'}^\top \phi_X)(a),
    \end{multline*}
    which implies $Y$ is a $T$-design. 
\end{proof}

\end{document}
